\clearpage
\section{Thực nghiệm bổ sung về chỉnh sửa đồ thị và dự đoán quan hệ}\label{app:probes}

\begin{table}[H]\centering\small
\caption{Các bước kiểm tra Task B. F1 là micro-F1 theo thao tác. Cỡ mẫu và
tập chấm được ghi theo từng phép thử do bộ nhãn thay đổi trong quá trình phát triển.}
\begin{tabular}{@{}p{.24\linewidth}p{.70\linewidth}@{}}\toprule
Phép thử & Kết quả và phạm vi diễn giải\\\midrule
Kiểm tra nhãn ảnh &
Người duyệt xác nhận 20/20 cặp train của bộ 437 cặp. Lượt kiểm tay trên
bộ 345 cặp ở thân bài dùng nhãn theo câu lệnh và lấy mẫu riêng.\\
REACT++ trên MagicBrush &
Trong 20 cặp validation, trợ lý đối chiếu câu lệnh với đỉnh, cạnh và chênh lệch giữa hai đồ thị.
Ba cặp thể hiện rõ thay đổi, sáu cặp thể hiện một phần và 11 cặp thiếu bằng chứng.
Ở lệnh đổi trụ cứu hỏa sang đỏ, đồ thị giữ được trụ nhưng thiếu thuộc tính màu.\\
Giải mã chuỗi thao tác &
Cả sáu lượt v1 đạt F1 bằng 0 và gặp lỗi giải mã trên cả 67 cặp validation.
Ở v2, MLP seed 42 dự đoán đúng một số thao tác nhưng tỷ lệ giao dịch hợp lệ
chưa đạt ngưỡng 50\%, nên hàng đợi dừng trước lượt GNN.\\
Học thuộc 16 mẫu &
GNN khớp toàn bộ đồ thị ở 15/16 mẫu train với nhãn đầy đủ và 16/16 sau khi rút
gọn thuộc tính. Hai lượt dùng cùng mẫu và seed 42, song đáp án thay đổi theo bộ nhãn.\\
Liên kết từ với đối tượng &
Trên 64 mẫu validation có thao tác sau khi rút gọn nhãn, thêm tín hiệu khớp từ
vào mô hình giúp F1 sau giải mã có ràng buộc tăng từ 11,14 lên 13,33\% với MLP
và từ 10,38 lên 10,90\% với GNN (seed 42).\\
Sao chép tên từ câu lệnh &
Trên 67 mẫu validation, sao chép giúp GNN gọi đúng tên ở 9/42 thao tác thêm đỉnh
so với 7/42 khi tắt sao chép. F1 toàn thao tác lần lượt là 10,07 và 10,10\%.
Cả bốn cấu hình MLP/GNN, bật/tắt sao chép, đều khớp toàn bộ đồ thị ở 0/67 mẫu (seed 42).\\
Thử hai bộ nhãn &
Cùng 14 mẫu train và seed 42, F1 train đạt 86,49\% với nhãn chênh lệch và
100\% với nhãn theo lệnh. Cả hai đạt 0\% trên 13 mẫu validation dùng chung.\\
Tái lập trên dữ liệu của He và cộng sự &
Qua sáu lượt MLP/GNN, cơ chế dừng sớm của v4 giảm tổng số epoch từ 72 xuống 61.
Các checkpoint được chọn cho dự đoán trùng với v3 trên 530 mẫu test.\\\bottomrule
\end{tabular}\end{table}

\begin{table}[H]\centering\small
\caption{Các đối chứng trước khi đánh giá PredCls trên tập test. MRR đo thứ hạng
của thực thể đích trong bài toán tri thức; mR@50 đo recall trung bình theo lớp quan hệ trong ảnh.}
\begin{tabular}{@{}p{.24\linewidth}p{.70\linewidth}@{}}\toprule
Phép thử & Kết quả trên validation\\\midrule
Quy tắc xếp hạng VG150 &
Nhóm chấm lại 54 checkpoint trên cùng 5.000 ảnh. Với GAT, mR@50 trung bình ba seed
là 19,55\% khi mỗi cặp giữ một predicate (GC), và 33,21\% khi mọi predicate
đều được tham gia xếp hạng toàn ảnh (NG).\\
Điểm quan hệ ConceptNet &
Ở nhánh fusion dùng điểm quan hệ của KG trên 5.000 ảnh VG150, mR@50 đạt
25,904\% với ánh xạ thật và 25,863\% khi tráo ánh xạ, trung bình ba seed theo
bộ chấm validation của lượt này.\\
Nhúng tri thức &
Trên 512 bộ ba validation của ConceptNet, MRR của TransE là 0,01567, ComplEx
là 0,01497 và RotatE v2 là 0,01324 (seed 42). Ở lượt RotatE v1, checkpoint
tốt nhất xuất hiện ngay tại epoch đầu.\\
Pilot GQA &
Với 512 ảnh train, 128 ảnh validation, seed 42 và năm epoch, GC mR@50 đạt
0,340\% với MLP và 0,315\% với GAT. Trong ba biến thể PE-Net, PE-CE đạt
0,215\%, PE-PL đạt 0,479\% và PE-PLPR đạt 0,441\%.\\\bottomrule
\end{tabular}\end{table}
